% Analise complementar por coleta. A validacao principal por leitura esta em
% ../validacao_por_leitura/secao_para_artigo.tex.
% Conferir no cadastro original a numeracao dos vasos
% inoculados: o CSV usa 1, 6, 7, 8 e 9, enquanto o metodo atual cita 6 a 10.
\subsection{Validação por coleta e por vaso}

As configurações de sensores selecionadas na análise combinatória foram
avaliadas com o mesmo classificador Extra Trees, sem novo ajuste de
hiperparâmetros. Cada coleta recebeu uma única previsão, obtida pela média
das probabilidades de suas leituras. A unidade de avaliação foi a coleta em
todos os protocolos. Para identificar as unidades biológicas, combinou-se
a classe experimental com o número do vaso, pois o número 1 aparece em ambos
os grupos no arquivo analisado.

Foram considerados três protocolos. O split original reservou 12 das 36
coletas para teste, aproximadamente um terço da base, mas oito dos dez vasos
aparecem tanto no treino quanto no teste. Em uma divisão adicional, sete vasos
foram destinados ao treino e três ao teste, com estratificação pela classe e
semente 42. Por fim, a validação \textit{leave-one-subject-out} (LOSO) realizou
dez rodadas; em cada uma, todas as coletas de um vaso foram testadas após
treinar o modelo nos outros nove vasos. A acurácia balanceada do LOSO foi
calculada após reunir as 36 previsões de coleta, pois cada rodada individual
contém apenas uma das classes.

\begin{table}[htbp]
\centering
\caption{Acurácia balanceada por coleta nos protocolos de validação. Os
valores entre parênteses indicam coletas corretamente classificadas.}
\label{tab:validacao_vaso_loso}
\begin{tabular}{lrrr}
\hline
Sensores & Split original & 70/30 por vaso & LOSO por vaso \\
\hline
MQ2, MQ135, MQ138 + ambiente & 91,67\% (11/12) & 91,67\% (9/10) & 89,16\% (32/36) \\
MQ2, MQ3, MQ8, MQ135, MQ138 & 91,67\% (11/12) & 91,67\% (9/10) & 80,65\% (29/36) \\
Seis MQ & 91,67\% (11/12) & 91,67\% (9/10) & 77,71\% (28/36) \\
Seis MQ + ambiente & 83,33\% (10/12) & 100,00\% (10/10) & 86,22\% (31/36) \\
MQ7 + ambiente & 75,00\% (9/12) & 91,67\% (9/10) & 86,53\% (31/36) \\
\hline
\end{tabular}
\end{table}

O conjunto MQ2, MQ135 e MQ138 acrescido das variáveis ambientais manteve o
maior resultado LOSO entre as cinco configurações avaliadas: 32 das 36
coletas corretas e 89,16\% de acurácia balanceada. No LOSO, acertou 16 das
17 coletas inoculadas e 16 das 19 coletas saudáveis. A avaliação 70/30 por
vaso envolve somente três vasos de teste e depende fortemente dos vasos
sorteados; por isso, o LOSO fornece uma visão mais abrangente dos dez vasos.
Como os conjuntos de sensores foram escolhidos anteriormente com estes mesmos
dados, os resultados constituem validação interna exploratória. A confirmação
final requer novos vasos ou campanhas independentes.
